% GENERATED FILE. Edit canonical lessons, then run npm run generate:books.
% canonical-source-hash: fnv1a64:7cf3aa3dc1268b0f
% canonical-lessons: ML-C552-bandham, ML-C552-dampatyam, ML-C552-pranayalekhanam, ML-C552-kamukan, ML-C552-pranayikkuka

\chapter{Love and Divorce}
\label{ch:ml-love-and-divorce}
\hlchaptermodality{552}

\begin{chapteropening}
\textbf{By the end of this chapter:} \emph{I can say a relationship, married life, a love letter, a boyfriend and to be in love.}

Complete the last lesson of chapter 552: I can say a relationship, married life, a love letter, a boyfriend and to be in love.
\end{chapteropening}

\section[\texorpdfstring{\ml{ബന്ധം} (bandhaṁ)}{ബന്ധം (bandhaṁ)}]{\ml{ബന്ധം} (bandhaṁ) — a relationship; a tie}

\label{lesson:ML-C552-bandham}



\begin{quote}
Before the new one: say the Malayalam for to put in work, then the Malayalam for to toil.
\end{quote}

\subsection*{You'll want to know: \ml{ബന്ധം}}
\textbf{\ml{ബന്ധം}} — \emph{bandhaṁ} — \textquotedblleft{}a relationship; a tie\textquotedblright{}.

\begin{quote}
\ml{സഹപ്രവർത്തകരുമായി} \ml{എനിക്ക്} \ml{നല്ല} \ml{ബന്ധമാണ്}.
\end{quote}

\begin{quote}
\emph{sahapravarttakarumāyi enikkŭ nalla bandhamāṇŭ.}
\end{quote}

\begin{quote}
I have a good relationship with my colleagues.
\end{quote}

\ml{ബന്ധു}, a relative, shares its root; \ml{ബന്ധപ്പെടുക} is to get in touch.

\begin{cousinweb}
\textbf{\ml{ബന്ധം}} (\emph{bandhaṁ}) is borrowed from Sanskrit \emph{bandha}, \textquotedblleft{}a bond, a tie\textquotedblright{} (CDIAL 9136).
\end{cousinweb}

\begin{grammarlens}[title={What you've built}]
One more word for describing what you lived and felt.
\end{grammarlens}

\subsection*{Your turn}
\begin{itemize}
\raggedright
  \item \emph{Say it:} \emph{bandhaṁ}
  \item \emph{Say it:} \emph{bandhaṁ}, once more
  \item \emph{Say it:} \emph{sahapravarttakarumāyi enikkŭ nalla bandhamāṇŭ.} — I have a good relationship with my colleagues.
  \item \emph{Recall:} say the Malayalam for to put in work, then the Malayalam for to toil, then say \emph{bandhaṁ} again
\end{itemize}

\begin{tcolorbox}[breakable,colback=teal!4,colframe=teal!35!black,title={Before you move on}]
What does \ml{ബന്ധം} mean? (\textquotedblleft{}A relationship; a tie\textquotedblright{}.) Say it once more.
\end{tcolorbox}
\section[\texorpdfstring{\ml{ദാമ്പത്യം} (dāmpatyaṁ)}{ദാമ്പത്യം (dāmpatyaṁ)}]{\ml{ദാമ്പത്യം} (dāmpatyaṁ) — married life}

\label{lesson:ML-C552-dampatyam}



\begin{quote}
Before the new one: say the Malayalam for to toil, then the Malayalam for a relationship.
\end{quote}

\subsection*{You'll want to know: \ml{ദാമ്പത്യം}}
\textbf{\ml{ദാമ്പത്യം}} — \emph{dāmpatyaṁ} — \textquotedblleft{}married life\textquotedblright{}.

\begin{quote}
\ml{മുത്തച്ഛന്റെയും} \ml{മുത്തശ്ശിയുടെയും} \ml{ദാമ്പത്യം} \ml{അമ്പത്} \ml{വർഷം} \ml{നീണ്ടു}.
\end{quote}

\begin{quote}
\emph{muttacchanṟeyuṁ muttaśśiyuṭeyuṁ dāmpatyaṁ ampatŭ varṣaṁ nīṇṭu.}
\end{quote}

\begin{quote}
My grandparents' married life lasted fifty years.
\end{quote}

\ml{ദമ്പതികൾ}, a married couple, is already familiar.

\begin{cousinweb}
\textbf{\ml{ദാമ്പത്യം}} (\emph{dāmpatyaṁ}) is borrowed from Sanskrit \emph{dāmpatya}, \textquotedblleft{}the state of husband and wife\textquotedblright{} (Monier-Williams), the noun of \emph{dampati}, \textquotedblleft{}a married couple\textquotedblright{}.
\end{cousinweb}

\begin{grammarlens}[title={What you've built}]
One more word for describing what you lived and felt.
\end{grammarlens}

\subsection*{Your turn}
\begin{itemize}
\raggedright
  \item \emph{Say it:} \emph{dāmpatyaṁ}
  \item \emph{Say it:} \emph{dāmpatyaṁ}, once more
  \item \emph{Say it:} \emph{muttacchanṟeyuṁ muttaśśiyuṭeyuṁ dāmpatyaṁ ampatŭ varṣaṁ nīṇṭu.} — My grandparents' married life lasted fifty years.
  \item \emph{Recall:} say the Malayalam for to toil, then the Malayalam for a relationship, then say \emph{dāmpatyaṁ} again
\end{itemize}

\begin{tcolorbox}[breakable,colback=teal!4,colframe=teal!35!black,title={Before you move on}]
What does \ml{ദാമ്പത്യം} mean? (\textquotedblleft{}Married life\textquotedblright{}.) Say it once more.
\end{tcolorbox}
\section[\texorpdfstring{\ml{പ്രണയലേഖനം} (praṇayalēkhanaṁ)}{പ്രണയലേഖനം (praṇayalēkhanaṁ)}]{\ml{പ്രണയലേഖനം} (praṇayalēkhanaṁ) — a love letter}

\label{lesson:ML-C552-pranayalekhanam}



\begin{quote}
Before the new one: say the Malayalam for a relationship, then the Malayalam for married life.
\end{quote}

\subsection*{You'll want to know: \ml{പ്രണയലേഖനം}}
\textbf{\ml{പ്രണയലേഖനം}} — \emph{praṇayalēkhanaṁ} — \textquotedblleft{}a love letter\textquotedblright{}.

\begin{quote}
\ml{അവൻ} \ml{അവൾക്ക്} \ml{ഒരു} \ml{പ്രണയലേഖനം} \ml{എഴുതി}.
\end{quote}

\begin{quote}
\emph{avan avaḷkkŭ oru praṇayalēkhanaṁ eḻuti.}
\end{quote}

\begin{quote}
He wrote her a love letter.
\end{quote}

\ml{പ്രണയം} love + \ml{ലേഖനം} a piece of writing.

\begin{cousinweb}
\textbf{\ml{പ്രണയലേഖനം}} (\emph{praṇayalēkhanaṁ}) is a learned compound built from Sanskrit \emph{praṇaya}, \textquotedblleft{}affection\textquotedblright{} (CDIAL 8525), and \emph{lekhana}, \textquotedblleft{}writing down\textquotedblright{} (Monier-Williams), a word CDIAL 11103 records in its older sense of scraping.
\end{cousinweb}

\begin{grammarlens}[title={What you've built}]
One more word for describing what you lived and felt.
\end{grammarlens}

\subsection*{Your turn}
\begin{itemize}
\raggedright
  \item \emph{Say it:} \emph{praṇayalēkhanaṁ}
  \item \emph{Say it:} \emph{praṇayalēkhanaṁ}, once more
  \item \emph{Say it:} \emph{avan avaḷkkŭ oru praṇayalēkhanaṁ eḻuti.} — He wrote her a love letter.
  \item \emph{Recall:} say the Malayalam for a relationship, then the Malayalam for married life, then say \emph{praṇayalēkhanaṁ} again
\end{itemize}

\begin{tcolorbox}[breakable,colback=teal!4,colframe=teal!35!black,title={Before you move on}]
What does \ml{പ്രണയലേഖനം} mean? (\textquotedblleft{}A love letter\textquotedblright{}.) Say it once more.
\end{tcolorbox}
\section[\texorpdfstring{\ml{കാമുകൻ} (kāmukan)}{കാമുകൻ (kāmukan)}]{\ml{കാമുകൻ} (kāmukan) — a boyfriend; a lover}

\label{lesson:ML-C552-kamukan}



\begin{quote}
Before the new one: say the Malayalam for married life, then the Malayalam for a love letter.
\end{quote}

\subsection*{You'll want to know: \ml{കാമുകൻ}}
\textbf{\ml{കാമുകൻ}} — \emph{kāmukan} — \textquotedblleft{}a boyfriend; a lover\textquotedblright{}.

\begin{quote}
\ml{അവൾ} \ml{കാമുകനോടൊപ്പം} \ml{സിനിമയ്ക്ക്} \ml{പോയി}.
\end{quote}

\begin{quote}
\emph{avaḷ kāmukanōṭoppaṁ sinimaykkŭ pōyi.}
\end{quote}

\begin{quote}
She went to the cinema with her boyfriend.
\end{quote}

A girlfriend is \ml{കാമുകി}.

\begin{cousinweb}
\textbf{\ml{കാമുകൻ}} (\emph{kāmukan}) is borrowed from Sanskrit \emph{kāmuka}, \textquotedblleft{}longing for, in love with\textquotedblright{} (Monier-Williams), with the Malayalam masculine ending -\ml{ൻ}.
\end{cousinweb}

\begin{grammarlens}[title={What you've built}]
One more word for describing what you lived and felt.
\end{grammarlens}

\subsection*{Your turn}
\begin{itemize}
\raggedright
  \item \emph{Say it:} \emph{kāmukan}
  \item \emph{Say it:} \emph{kāmukan}, once more
  \item \emph{Say it:} \emph{avaḷ kāmukanōṭoppaṁ sinimaykkŭ pōyi.} — She went to the cinema with her boyfriend.
  \item \emph{Recall:} say the Malayalam for married life, then the Malayalam for a love letter, then say \emph{kāmukan} again
\end{itemize}

\begin{tcolorbox}[breakable,colback=teal!4,colframe=teal!35!black,title={Before you move on}]
What does \ml{കാമുകൻ} mean? (\textquotedblleft{}A boyfriend; a lover\textquotedblright{}.) Say it once more.
\end{tcolorbox}
\section[\texorpdfstring{\ml{പ്രണയിക്കുക} (praṇayikkuka)}{പ്രണയിക്കുക (praṇayikkuka)}]{\ml{പ്രണയിക്കുക} (praṇayikkuka) — to be in love}

\label{lesson:ML-C552-pranayikkuka}



\begin{quote}
Before the new one: say the Malayalam for a love letter, then the Malayalam for a boyfriend.
\end{quote}

\subsection*{You'll want to know: \ml{പ്രണയിക്കുക}}
\textbf{\ml{പ്രണയിക്കുക}} — \emph{praṇayikkuka} — \textquotedblleft{}to be in love\textquotedblright{}.

\begin{quote}
\ml{കോളേജിൽ} \ml{വെച്ച്} \ml{അവർ} \ml{പ്രണയിച്ചു}.
\end{quote}

\begin{quote}
\emph{kōḷējil veccŭ avar praṇayiccu.}
\end{quote}

\begin{quote}
They fell in love at college.
\end{quote}

Past: \ml{പ്രണയിച്ചു}, \emph{praṇayiccu}.

The noun is \ml{പ്രണയം}, romantic love; \ml{പ്രണയവിവാഹം} is a love marriage.

\begin{cousinweb}
\textbf{\ml{പ്രണയിക്കുക}} (\emph{praṇayikkuka}) is built from Sanskrit \emph{praṇaya}, \textquotedblleft{}guidance; affection\textquotedblright{} (CDIAL 8525), with the Malayalam verb ending -\ml{ിക്കുക}.
\end{cousinweb}

\begin{grammarlens}[title={What you've built}]
One more word for describing what you lived and felt.
\end{grammarlens}

\subsection*{Your turn}
\begin{itemize}
\raggedright
  \item \emph{Say it:} \emph{praṇayikkuka}
  \item \emph{Say it:} \emph{praṇayikkuka}, once more
  \item \emph{Say it:} \emph{kōḷējil veccŭ avar praṇayiccu.} — They fell in love at college.
  \item \emph{Recall:} say the Malayalam for a love letter, then the Malayalam for a boyfriend, then say \emph{praṇayikkuka} again
\end{itemize}

\begin{tcolorbox}[breakable,colback=teal!4,colframe=teal!35!black,title={Before you move on}]
What does \ml{പ്രണയിക്കുക} mean? (\textquotedblleft{}To be in love\textquotedblright{}.) Say it once more.
\end{tcolorbox}
